% ============================================================================
%  Resumen
% ============================================================================
\section*{Resumen}
\addcontentsline{toc}{section}{Resumen}

Se diseñó e implementó el control de una plataforma móvil diferencial de dos
motorreductores DC con un PIC16F877A y el doble puente H DRV8833. La velocidad
de cada rueda se controla por PWM de hardware (CCP1 y CCP2 sincronizados con
Timer2) a \textbf{\SI{1.25}{\kilo\hertz} con 1000 niveles de ciclo de trabajo
($\approx$10~bits)}, y el sentido de giro y el frenado dinámico se gobiernan
con dos líneas digitales de dirección.

El firmware, escrito en C con MPLAB XC8, está organizado en tres capas
(hardware, control por motor y navegación) y ofrece dos modos seleccionables en
compilación: la secuencia automática de validación del taller (avance, reversa,
giro y frenado con rampas) y el control remoto desde un celular por Bluetooth
(HC-05/HC-06, UART a 9600~baudios, con interrupción de recepción y control de
velocidad en 11 niveles).

Durante el diseño se identificó que la frecuencia de \SI{1.2}{\kilo\hertz}
pedida no es alcanzable con cristal de \SI{20}{\mega\hertz} y prescaler 1:16
(exige $PR2 = 259 > 255$); se justifica el uso de \SI{1.25}{\kilo\hertz}.
También se resolvieron la ausencia del pulso de habilitación nSLEEP y del giro
en el código base, y la diferencia entre el cristal de la lista de materiales
(\SI{4}{\mega\hertz}) y el del código (\SI{20}{\mega\hertz}).

\medskip
\noindent\textbf{Palabras clave:} PIC16F877A, DRV8833, PWM, Timer2, robótica
móvil diferencial, puente H, UART, Bluetooth, MPLAB XC8, sistemas embebidos.
